Many NoSQL stores lack multi-key transactions, so distributed transactions spanning several of them
are commonly executed through a shim that runs each branch under strict two-phase locking (S2PL)
and commits it with two-phase commit (2PC). Under 2PC, however, a branch's writes are final as soon
as it finishes executing, while S2PL holds its locks through the prepare phase, the coordinator's
round trips and the commit, so conflicting transactions wait for values that will no longer change.
We present a speculative shim that lets a waiting transaction read these final but uncommitted
values instead. It enforces commitment ordering, which preserves global serializability across
shims without any global concurrency control. The cost is cascading aborts: when a transaction
aborts, every transaction that speculated on it aborts too, and under high contention this cost
outgrows the waiting saved. To combine the strengths of both approaches, an adaptive variant caps
how many transactions may speculate on a single uncommitted writer and tunes this cap at run time
from the shim's commit rate, moving between full speculation and the non-speculative shim. On a
micro-benchmark, speculation reaches a 10--100\% higher peak throughput than the non-speculative
shim, with 2--8$\times$ lower latency, and gains most with skewed access and long-running
transactions. When many branches of a transaction are skewed, the advantage disappears, since each
skewed branch is another independent chance of a cascading abort; when only one branch is skewed,
speculation keeps winning, though the gap narrows as branches are added. Uncapped speculation
degrades past its peak as the number of threads grows, whereas the adaptive shim degrades less and
follows the better of the two shims across thread counts.

\vspace{1em}
\noindent\textbf{Keywords:} distributed transactions, NoSQL, two-phase commit, commitment ordering,
speculative concurrency control, cascading aborts
